% \section{Challenges and Future Directions}

% The evidence reviewed in this survey characterizes current RSI as a collection of bounded improvement loops rather than a fully autonomous recursive process. The three bottlenecks identified in the introduction remain only partially resolved: (1) foundation-model development requires substantial resources and coordination; (2) scalable learning depends heavily on externally constructed experience and verification; and (3) adaptation after deployment demands recurring human effort to diagnose, validate, and release updates. We identify five potential research directions toward genuine RSI. %  Addressing these bottlenecks requires systems that can identify what limits progress, acquire useful experience, and carry validated improvements into subsequent rounds. Beyond automating these activities, genuine RSI requires the resulting capabilities to strengthen the process through which further improvements are discovered and retained.

% \noindent\textbf{Cross-Component Improvement.}
% AI capabilities emerge from interactions among data, model architectures, training methods, infrastructure, and deployment environments. Improving one component in isolation may yield diminishing returns or shift the bottleneck elsewhere. Future RSI systems should use experimental evidence to locate the constraints on progress, distinguish their causes, and coordinate changes across components. For example, repeated task failures may call for additional training data, a revised learning objective, or a better tool interface; choosing among these interventions requires estimating their expected benefits and costs. A central research challenge is to attribute gains when several components change together and to allocate limited resources as bottlenecks shift. This would move RSI from optimizing predefined targets toward autonomously organizing the improvement process.

% \noindent\textbf{Adaptive Acquisition of Learning Experience.}
% The value of an experience depends on what the learner can already do and what it is ready to learn next. Future RSI systems should diagnose capability gaps and actively select, generate, or seek experiences that address them. Relevant mechanisms include constructing tasks near the learner's capability frontier, redesigning environments to expose informative failures, and requesting external feedback when internal verification is insufficient. The challenge is to distinguish experiences that produce transferable learning from those that merely increase apparent difficulty or reinforce existing errors. Acquisition strategies should therefore be evaluated by the durable capability gains they enable relative to the cost of generating, verifying, and learning from experience. As the learner improves, these strategies should adapt, allowing experience acquisition and capability development to advance together.

% \noindent\textbf{Governed Continual Improvement after Deployment.}
% Deployed systems must turn operational feedback into persistent improvements while maintaining reliable service. Existing practices partially automate this process: Meta's agents, for example, can investigate performance regressions and generate pull requests, while engineers retain responsibility for reviewing proposed changes~\cite{meta_capacity_efficiency_2026}. Future research should develop mechanisms that distinguish isolated failures from recurring limitations, determine when an update is warranted, and select an appropriate scope of intervention. Proposed changes should undergo controlled validation, followed by monitoring to establish whether their benefits persist under changing requirements and environments. Reversible checkpoints, sandboxed execution, and staged releases can support this process, with human review retained for consequential changes. When an update fails to generalize, the system should revise its diagnosis and improvement strategy, making failed interventions informative for subsequent rounds.

% \noindent\textbf{Trustworthy Evolution of Improvement Mechanisms.}
% A defining step toward genuine RSI is to make the mechanisms that diagnose, generate, evaluate, and select improvements themselves subject to improvement. A system might learn better search strategies, revise its failure analysis, or develop more informative validation procedures. However, changing these mechanisms also changes how successor systems are produced and judged, creating opportunities for evaluator drift, reward exploitation, and the loss of effective search strategies. Future work should investigate how to validate such revisions against protected external criteria and preserve the provenance of inherited changes. The central question is whether a revised mechanism enables more effective subsequent improvement under comparable constraints, beyond any immediate gain in task performance.

% \noindent\textbf{Long-Horizon Recursive Evaluation.}
% Single-round gains and final benchmark scores provide limited evidence of recursive progress. Future evaluations should track sequences of inherited updates, testing whether improvements remain useful, transfer across tasks, preserve earlier capabilities, and increase the effectiveness of later improvement rounds. Comparisons should account for compute, data, tool use, and human intervention, and distinguish gains from inherited changes from those obtained through additional search or repeated retraining. Protected test sets and changing task streams can help reveal regressions, overfitting, and evaluator drift. Such protocols should establish when improvement compounds, when it plateaus, and when continued adaptation becomes too costly or unreliable.


% \iffalse
% \section{Challenges and Future Directions}

% The evidence reviewed in this survey suggests that current RSI is better characterized as a collection of bounded improvement loops than as a fully autonomous recursive process.
% % Existing systems can execute prescribed updates, search over prompts and agent implementations, generate learner-conditioned experience, adapt persistent components from deployment feedback, and in a smaller number of cases revise mechanisms that govern later improvements.
% However, the three bottlenecks motivating RSI at the beginning of this survey remain only partially resolved: (1) foundation-model development is still expensive to coordinate, (2) scalable learning still depends heavily on externally constructed experience and verification, and (3) deployed systems still require substantial human effort to diagnose, validate, and release persistent updates. Progress toward genuine RSI therefore requires not merely increasing the number of components that AI may modify, but improving the reliability, efficiency, and inheritance of the entire improvement process.
% % We identify five directions that appear particularly important: cross-component improvement, autonomous acquisition of useful experience, long-horizon recursive evaluation, trustworthy evolution of improvement mechanisms, and governed adaptation in open environments. Together, these directions extend the autonomy hierarchy introduced in this survey while preserving the efficiency and innovation objectives that motivate RSI.

% \noindent\textbf{Cross-Component Improvement.}
% Future systems should infer from experimental evidence which components currently constrain progress and dynamically reallocate improvement effort as these bottlenecks shift.

% \noindent\textbf{Adaptive Acquisition of Learning Experience.}
% Future RSI should identify and acquire the experience with the greatest expected learning value at the learner's current capability frontier, rather than relying on fixed or externally constructed experience distributions.

% \noindent\textbf{Governed Real-World RSI}
% Future studies should evaluate persistent updates across long task streams in which requirements and environments change over time. They should test whether each update is invoked when relevant, continues to improve later tasks, and preserves capabilities acquired earlier. Protected test sets keep part of the evaluation outside the adaptive loop and can expose both regressions and evaluator drift. Reversible checkpoints and sandboxed execution allow faulty changes to be isolated and withdrawn, while consequential releases should still undergo human review.

% \noindent\textbf{Controlled Continual Improvement in Real-World Environments.}
% Deployed RSI systems need to translate problems encountered during operation into lasting system improvements. Existing engineering practices have partially automated this process. For example, Meta's agents can investigate performance regressions and generate pull requests, although the proposed changes still require review by engineers~\cite{meta_capacity_efficiency_2026}. Future research should explore continual update mechanisms driven by operational feedback, enabling systems to distinguish isolated failures from persistent issues that require system updates and determine when an improvement is warranted. A key challenge is that a fix for the current case may not generalize to subsequent tasks, while gains observed in offline testing may not persist in changing environments. Research is therefore needed on how to validate the effects of proposed changes under controlled conditions and use evidence from subsequent operation to decide whether to retain them. When improvements fail to deliver sustained benefits, systems should be able to reassess the underlying problems and adjust their strategies, reducing the engineering effort required to review and deploy individual updates while maintaining service stability.


% \noindent\textbf{Trustworthy Evolution of Improvement Mechanisms.}
% Future RSI should make the mechanisms that diagnose, generate, evaluate, and select improvements themselves subject to revision, while preserving protected external anchors that distinguish genuine capability gains from changes in the criteria by which successors are judged.

% \noindent\textbf{Long-Horizon Recursive Evaluation.}
% Future RSI should be evaluated over sequences of inherited improvements rather than by a single end-point score, testing whether gains persist, transfer, and compound across generations under matched resource budgets and protected evaluation.

% \noindent\textbf{Toward Genuine Recursive Self-Improvement.}
% Taken together, these directions suggest a transition from today’s bounded improvement systems toward genuine RSI.
% The next generation of systems will need to decide not only how to execute an improvement, but also what should be improved, what experience should be acquired, which mechanisms should govern subsequent improvement rounds, and whether continued improvement is justified by its expected gain, cost, and risk.
% % At the same time, each accepted update must remain attributable, externally testable, and reversible.

% \fi


%\clearpage
\section{Challenges and Future Directions}

The preceding chapters show progress in automating improvement execution, strategy selection, experience acquisition, deployment adaptation, and the revision of improvement mechanisms. They also establish that greater autonomy, durable capability gains, and a more effective improvement process are distinct properties. The three bottlenecks motivating this survey therefore remain only partially resolved: foundation-model development requires substantial resources and coordination; scalable learning depends on useful experience and reliable verification; and deployed systems require continuing effort to diagnose, validate, and release updates. Building on the autonomy hierarchy, the four application regimes, and the industrial practices reviewed above, we identify eight research directions for connecting these partial advances into sustained recursive improvement.

\noindent\textbf{Cross-Component Diagnosis and Coordinated Improvement.}
The transition from L1 to L2 exposes a diagnostic problem: an observed failure rarely identifies the component that should change. An incorrect answer may originate in defective source data, missing context, a tool interface, the model, or the evaluator. The Humanlaya case in Section~\ref{subsec:humanlaya} makes this ambiguity concrete: a failed extraction was initially treated as evidence of poor document quality, while the appropriate intervention concerned the quality checker's decision rule. In science and embodied intelligence, the same problem extends to experimental protocols, perception, control, and environmental conditions. As several components adapt, changing one can also invalidate assumptions on which another relies.

Future systems should turn diagnoses into testable intervention hypotheses and use controlled comparisons to estimate individual effects and interactions. Selective component freezing, targeted ablations, and explicit records of dependencies could help distinguish a useful repair from compensation for an upstream defect. Coordinated search should then allocate effort across data, models, harnesses, and environments as constraints shift. Theseus's workspace studies in Section~\ref{sec:practice_theseus} illustrate the importance of separating environment improvements from changes to the task model; the reported component-level gains leave the full learning cycle to be established. The research goal is to identify which intervention improves the overall system, under what conditions, and whether the resulting knowledge guides later improvement decisions.

\noindent\textbf{Learner-Conditioned Experience Acquisition and Reliable Learning Signals.}
The L3 analysis distinguishes three properties of experience: whether it is valid, how difficult it is for the current learner, and whether learning from it produces a durable benefit. AZR uses executable checks alongside learner-dependent task proposal, whereas R-Zero uses solver consistency as a proxy for difficulty~\cite{zhao2025absolutezero,huang2026rzero}. Neither correctness nor estimated difficulty alone establishes learning value. An adaptive curriculum may repeatedly select ambiguous tasks, reinforce evaluator errors, or concentrate on a narrow region where progress is easy to measure. Such errors can enter both the learner and the state that determines its next experiences.

Future acquisition mechanisms should estimate learning value across multiple rounds while preserving validity, diversity, and coverage of previously acquired capabilities. This includes selecting from existing data, generating new tasks, seeking interactions, and requesting stronger external feedback when internal judgments are insufficient. A central challenge is deciding when expensive verification is worth its cost and how uncertain experience should influence subsequent training. Evaluation should allow the acquired distributions to differ while matching data-source access and total acquisition and learning budgets. Freezing the learner-state input to the acquisition mechanism, or replacing adaptive selection with a learner-independent schedule, can test whether learner conditioning improves transfer beyond the effects of additional training.

\noindent\textbf{Persistent-State Management and Reliable Reuse.}
L4 systems inherit parameters, memories, skills, tools, and harness code, but retaining an artifact does not establish that later agents benefit from it. The L4 analysis identifies failures in both activation and execution: a relevant skill may never be retrieved, or an agent may retrieve it without following it correctly. Accumulation creates an additional difficulty. Library Drift shows that an expanding skill library can degrade retrieval and stall improvement, while overly aggressive retirement can also be harmful~\cite{zhang2026library}. Persistence therefore requires managing the continuing applicability and interaction of retained changes.

Future work should associate inherited artifacts with their supporting evidence, applicability conditions, dependencies, and observed effects on later tasks. Admission tests such as HDSO's paired evaluations provide a starting point~\cite{shang2026hypothesis}, but validation must continue as the task distribution, executor, and surrounding system change. Studies should separately measure update quality, activation, faithful use, and downstream benefit, then investigate when to merge, revise, retire, or revalidate artifacts. Cross-model reuse introduces a further question: a skill useful to its author may be unsuitable for a different executor. Resolving these issues would connect successful update generation to reliable capability transfer across sessions and successors.

\noindent\textbf{Governed Adaptation under Domain-Specific Feedback.}
The four application regimes show why a common improvement loop requires different validation strategies. In science, an unsuccessful experiment may reflect the hypothesis, protocol, or instrument, and evidence must retain its assumptions and uncertainty. In embodied intelligence, the current policy determines which states are observed, while physical trials consume resources and can have irreversible consequences. Software offers executable feedback, but passing tests remains conditional on the specification and test coverage. In healthcare, outcomes are delayed and confounded, and an update validated in one population or institution may not transfer to another. These differences constrain both what a system can learn autonomously and what evidence supports deployment.

Future research should develop update policies that distinguish transient failures from recurring limitations and select a validation process appropriate to the proposed change. Simulation, retrospective replay, and sandboxed execution can support initial screening, followed by supervised or staged evaluation under the intended operating conditions. Monitoring should test whether benefits persist as requirements, populations, tools, and environments change. Restoring a software checkpoint cannot undo physical or clinical consequences, making pre-release validation essential in those settings. Human authority over consequential releases, access permissions, and safety constraints should remain explicit. The open problem is how to increase useful adaptation while controlling the propagation of errors across future interactions.

\noindent\textbf{Trustworthy Evolution of Improvement Mechanisms.}
At L5, changes to the improver, evaluator, or research policy alter how subsequent successors are produced and accepted. This creates a coupled validation problem: a stronger solver can expose weaknesses in its evaluator, but changing the evaluator can also make scores incomparable or reward exploitation. RQGM addresses part of this problem by freezing its evaluator within an epoch, validating replacements against an independent anchor, and revisiting scores affected by replacement~\cite{iacob2026redqueen}. The evaluator refinement described in Hyra raises the same broader question of how adaptive measurement can remain credible as search progresses.

Future systems should distinguish editable internal feedback mechanisms from independently maintained acceptance criteria and preserve evidence linking each mechanism revision to later decisions. Controlled comparisons should assess changes to candidate generation and evaluation separately before attributing gains to their joint evolution. A-Evolve-Training provides an example of policy inheritance within a fixed high-level objective~\cite{shi2026aevolvetraining}; whether such policies transfer beyond the research setting in which they evolved remains to be tested. Open questions include how to maintain independent assessment under repeated adaptive access, detect coordinated proposer--evaluator errors, and reverse harmful mechanism changes without losing useful experience. Revising operational research priorities need not transfer authority over the system's overall mission.

\noindent\textbf{Long-Horizon Evaluation of Inherited Improvement Capacity.}
The distinction between structural and effective L5 should guide evaluation. Structural evidence establishes that a revised mechanism is inherited and invoked; effectiveness requires showing that it produces or selects better subsequent improvements. This remains a substantive empirical gap. For example, the AIDE$^{2}$ experiment reviewed in this survey did not establish a statistically significant efficiency advantage when an evolved harness was installed as the outer improver~\cite{weco2026aide2}. Software lineage studies further motivate separating current task performance from the ability to produce useful descendants, since temporarily weaker variants may enable later progress.

Future benchmarks should compare original and revised mechanisms from comparable initial systems and evidence under matched total budgets, including the cost of developing and evaluating the mechanisms. Holding an evolved mechanism fixed on fresh tasks can test whether its improvement strategy transfers, as illustrated by HyperAgents~\cite{zhang2026hyperagents}. Longer studies should report complete trajectories across repeated runs, including rejected updates, regressions, recovery, retained capabilities, and resource use. Protected reference tasks can preserve comparability, while fresh tasks and unfamiliar constraints test transfer and discovery. These protocols should establish whether inherited changes alter subsequent improvement capacity, when gains plateau, and whether apparent acceleration survives accounting for additional search and evaluation effort.

\noindent\textbf{Resource-Aware Improvement and Human Collaboration.}
The efficiency objective motivating RSI requires accounting for the entire improvement process. A faster kernel, a better training recipe, or a longer unattended run may still require substantial candidate generation, evaluation, infrastructure maintenance, and human review. Meta's regression-analysis workflow retains engineering review of proposed changes~\cite{meta_capacity_efficiency_2026}; Humanlaya also retains human validation and reports handling time alongside quality. These practices show why autonomy should be assessed together with the amount and type of human effort required. Stronger base models and more parallel trials can further confound comparisons between improvement strategies.

Future work should study how to allocate a total budget among diagnosis, experience acquisition, candidate search, verification, and deployment monitoring. Cheap screening can reduce expensive trials, provided its predictions remain calibrated against independent outcomes. Reusing validated tools and prior failures may reduce repeated work, but its maintenance cost must also be counted. Evaluations should report time and total cost to a validated capability target, alongside review effort and rework. Adaptive stopping is another research priority: systems should justify continued experimentation by its expected learning value and uncertainty, with explicit conditions for pausing, escalating to a human, or terminating an unproductive search.

\noindent\textbf{Reproducible Infrastructure for Cross-Round Inheritance.}
Successor systems need access to how an improvement was obtained, including unsuccessful alternatives and the conditions under which its evidence holds. The industrial cases expose complementary requirements: Lark emphasizes temporally grounded data and failure attribution; Humanlaya retains versioned quality-system updates; Hyra stores executable experience; and Agent-Native Research Lab proposes research artifacts that preserve code, specifications, exploration history, and empirical traces. These practices motivate a shared infrastructure for carrying research evidence across rounds, while the full longitudinal benefit of such inheritance remains to be established.

A reusable artifact format should connect the parent state, proposed change, motivating evidence, evaluation configuration, acceptance decision, and subsequent use. Versioned data, environments, and evaluators would make results easier to replay and clarify which comparisons remain valid after an update. Deterministic re-extraction can check reported measurements, and formal verification can establish properties relative to a specification; neither alone establishes that the experiment or specification captures the intended capability. Public evaluation should therefore pair inspectable artifacts with replay studies and report result provenance, artifact availability, and independent replication separately. Where data are confidential, shareable evaluation subsets and controlled replay interfaces could support partial verification while preserving the limits of what can be checked.

Across these directions, progress toward genuine RSI should be demonstrated through repeated, attributable, and transferable improvements in the capacity to improve. The decisive evidence is that inherited changes help later rounds acquire more useful experience, discover better interventions, or validate successors more effectively under explicit resource and authority constraints.